\section{Released positive-viscosity profile: gauge-invariant curvature measures and the quantum-state map}
\label{sec:released-profile-measures}

The preceding non-Abelian embedding gives pointwise axis information.  The released positive-viscosity source supplies the stronger spatial estimates needed to pass from an axis value to an actual gauge-invariant spatial integral.  We use its original coordinates, viscosity and terminal variable
\[
  \tau=1-s,\qquad \nu>0,\qquad c>0,\qquad g>0,
\]
and retain the fixed nonzero bridge coefficient $\lambda$.  The source reader and its endpoint continuation prove that, for positive constants $C_\Omega$ and $C_{\dot u}$ fixed by the source profile,
\[
 \|\operatorname{curl}u_\nu(1-\tau)\|_2^2
 \ge \nu^{3/2}C_\Omega\tau^{-1/2-3h},\qquad
 \|\partial_su_\nu(1-\tau)\|_2^2
 \ge \nu^{5/2}C_{\dot u}\tau^{-3/2-3h},
 \tag{250}
\]
for the source range $0<h<1/100$ and sufficiently small $\tau>0$.  These are lower bounds on the full spatial integrals, not evaluations on the measure-zero axis.

Use the exact reducible connection in the original Cartesian chart,
\[
 A_0=0,\qquad A_i=\lambda u_{\nu,i}T,
 \qquad T=-i\sigma_3/2,
 \qquad -2\operatorname{tr}(T^2)=1.
 \tag{251}
\]
Since every component has the same Lie generator, $[A_\mu,A_\nu]=0$ and
\[
 F_{0i}=\lambda c^{-1}\partial_su_{\nu,i}T,
 \qquad
 F_{ij}=\lambda(\partial_i u_{\nu,j}-\partial_j u_{\nu,i})T.
 \tag{252}
\]
The two positive component norms are therefore exactly
\[
\begin{aligned}
 \mathcal K_B(s)&:=-2\int\sum_{i<j}\operatorname{tr}(F_{ij}^2)\,dx
      =\lambda^2\|\operatorname{curl}u_\nu(s)\|_2^2,\\
 \mathcal K_E(s)&:=-2\int\sum_i\operatorname{tr}(F_{0i}^2)\,dx
      =\lambda^2c^{-2}\|\partial_su_\nu(s)\|_2^2.
\end{aligned}
\tag{253}
\]
Combining (250) and (253) gives the exact terminal lower bounds
\[
 \mathcal K_B(1-\tau)\ge\lambda^2\nu^{3/2}C_\Omega\tau^{-1/2-3h},\qquad
 \mathcal K_E(1-\tau)\ge\lambda^2c^{-2}\nu^{5/2}C_{\dot u}\tau^{-3/2-3h}.
 \tag{254}
\]
Hence both component norms diverge as $\tau\downarrow0$, while the source energy identity gives
\[
 \int_0^1\mathcal K_B(s)\,ds<\infty,
 \qquad
 \int_0^1\mathcal K_E(s)\,ds=\infty.
 \tag{255}
\]
The first statement uses the original positive-viscosity dissipation bound; the second follows by integrating the nonintegrable exponent in (254).  These are Euclidean component norms.  They do not assert a sign for the Lorentzian action, whose metric has signature $(-,+,+,+)$.

The full sourced Yang--Mills equation remains coordinate-level precise.  With $D^\mu F_{\mu\nu}=g^2j_\nu$,
\[
 j_0=0,
 \qquad
 j_i=\frac{\lambda}{g^2}
 \left(\Delta u_{\nu,i}-c^{-2}\partial_s^2u_{\nu,i}\right)T,
 \tag{256}
\]
and substitution of the complete positive-viscosity equation gives
\[
 j_i=\frac{\lambda}{g^2}\left[\Delta u_{\nu,i}
 -c^{-2}\left(\nu\Delta(\partial_su_{\nu,i})
 -(\partial_su_\nu\!\cdot\!\nabla)u_{\nu,i}
 -(u_\nu\!\cdot\!\nabla)\partial_su_{\nu,i}
 -\partial_i\partial_sp_\nu+\partial_sf_{\nu,i}\right)\right]T.
 \tag{257}
\]
Every viscosity, pressure, force and transport term is displayed.  Taking the spatial divergence of (256) gives $D^\nu j_\nu=0$ from incompressibility, so the sourced map is internally consistent but is not a source-free Yang--Mills solution.

At each $s<1$, the source profile is smooth and compactly supported.  The input to the spatial observable map is the pointwise gauge-invariant density, not the already integrated scalar.  Define
\[
\begin{aligned}
 h_s^{B}(x)&:=-2\sum_{i<j}\operatorname{tr}(F_{ij}(s,x)^2)\\
            &=\lambda^2|\operatorname{curl}u_\nu(s,x)|^2,
\end{aligned}
\]
\[
\begin{aligned}
 h_s^{E}(x)&:=-2\sum_i\operatorname{tr}(F_{0i}(s,x)^2)\\
            &=\lambda^2c^{-2}|\partial_su_\nu(s,x)|^2.
\end{aligned}
\]
\[
 h_s:=h_s^{B}\quad\text{or}\quad h_s:=h_s^{B}+h_s^{E},\qquad
 H_{h_s}:=\int_{\mathbb R^3}h_s(x)\,dx.
\]
Thus $H_{h_s}=\mathcal K_B(s)$ for the magnetic choice and $H_{h_s}=\mathcal K_B(s)+\mathcal K_E(s)$ for the total choice.  For each lattice spacing $a>0$ and box size $L$, write an oriented edge as $e=(n,i)$ and define its physical midpoint $x_e=a(n+\tfrac12e_i)$.  The exact sampling morphism and face extension are
\[
 I_{L,a}:C_c^\infty(\mathbb R^3)\longrightarrow\mathbb R^{\mathsf E_L},\qquad (I_{L,a}h)_e=h(x_e),\qquad
 (I^P_{L,a}h)_p=\frac14\sum_{e\in\partial p}(I_{L,a}h)_e.
\]
It sends the pointwise continuum density to the original edge and face coefficients, with no averaging or normalization of $h$.  The same map has an explicit physical-volume comparison.  If $h\in C_c^1(\mathbb R^3)$ is supported in a compact set $K$ and the box contains all edge midpoints meeting $K_{a\sqrt3/2}$, then for each direction $i$ the midpoint cubes form a partition and
\[
 \left|a^3\sum_{e=(n,i)}h(x_e)-\int_{\mathbb R^3}h(x)\,dx\right|
 \le \frac{\sqrt3}{2}a\,\operatorname{vol}(K_{a\sqrt3/2})\,\|\nabla h\|_\infty.
\]
Indeed, integrate $h(x_e)-h(x)$ over the cube centered at $x_e$ and use the mean-value bound $|h(x_e)-h(x)|\le(\sqrt3a/2)\|\nabla h\|_\infty$; cubes outside $K_{a\sqrt3/2}$ contribute zero.  Thus the exact midpoint morphism preserves the unnormalized curvature mass in the joint lattice limit, with a displayed $O(a)$ error for each fixed profile.  In the following display, $D_{h,g}$ abbreviates this $D_{L,a,g}[h]$ at the displayed fixed regulator.  The resulting operator is exactly
\[
 D_{L,a,g}[h]=\kappa\sum_{e\in\mathsf E_L}(I_{L,a}h)_eE_e
 +b\sum_{p\in\mathsf P_L}(I^P_{L,a}h)_p(2-W_p),
\]
so the full edge/face sets and the physical scales remain present.  This is the domain-correct composition $h\mapsto I_{L,a}h\mapsto D_{L,a,g}[h]\mapsto\Xi_{L,a,g}[h]$ used below.  Apply the already proved exact observable map
\[
 h_s\longmapsto D_{h_s,g}
 \longmapsto \Xi_{h_s,g}=
 (D_{h_s,g}-\langle\psi_g,D_{h_s,g}\psi_g\rangle)\psi_g
 \longmapsto e^{-tA_g/2}\Xi_{h_s,g}.
 \tag{258}
\]
The centering is exact and no division by the raw norm occurs.  Gauge invariance is a projection, not a normalization: if $\mathcal G_L$ is the original vertex gauge group with product Haar probability, then
\[
 (\Pi_G f)(U)=\int_{\mathcal G_L}f(g\mathbin{\cdot}U)\,dg,
 \qquad \Pi_G^*=\Pi_G,\quad \Pi_G^2=\Pi_G,\quad \|\Pi_Gf\|\le\|f\|.
\]
Haar invariance proves the two identities and the contraction inequality.  The operators $D_{L,a,g}[h]$ and the vacuum commute with this action, so $\Pi_G\Xi_{L,a,g}[h]=\Xi_{L,a,g}[h]$ exactly; no scalar rescaling is introduced.  Vacuum centering removes only the component along $\psi_g$, whereas division by a state norm would define a separate normalized spectral probability and is not performed here.  The common free spectral map gives, for each fixed $s<1$ and $t>0$,
\[
 \|J_th_s\|^2\le \frac{9}{\pi^4t^8}\,\|h_s\|_1^2,
 \qquad
 H_{h_s}=\int h_s\,dx=\lambda^2\|\operatorname{curl}u_\nu(s)\|_2^2>0
\tag{259}
\]
whenever the profile is nonzero.  Thus (254) proves divergence of the input mass $H_{h_s}$ at the terminal endpoint, but (259) also shows why that divergence alone does not prove a zero interacting mass gap: the centered quantum vector, its actual interacting norm, and the joint regulator/volume limit still require estimates not supplied by the classical source.  The exact propagation established here is the map (251)--(258), with all source terms and units retained.

The inspected source provenance is
\texttt{ns\_inner\_profile\_curvature\_current\_bridge.md}.
The released reader and endpoint/axis continuation are recorded in the lane
provenance files.  No theorem about a source-free four-dimensional Yang--Mills
mass gap is inferred from the sourced profile.



\section{A support-uniform lower bound for the free curvature state}
\label{sec:support-uniform-free-state}

The compact support statement of the released source yields a quantitative consequence that uses the unnormalized state directly.  Let $K\subset\mathbb R^3$ be a fixed compact support for $h_s\ge0$, and choose $R>0$ with $K\subset\{x:|x|\le R\}$.  Write $H_s=\int h_s(x)\,dx$.  For every $k$ with $|k|\le(2R)^{-1}$,
\[
 |\widehat h_s(k)-H_s|
 =\left|\int_K(e^{-ik\cdot x}-1)h_s(x)\,dx\right|
 \le |k|R H_s\le \tfrac12H_s,
\]
so $|\widehat h_s(k)|\ge H_s/2$.  Inserting this exact bound into the already proved free norm formula gives
\[
 \|J_t h_s\|^2
 \ge \frac{H_s^2}{1280\pi^5t}
       4\pi\int_0^{(2R)^{-1}}r^6e^{-tr}\,dr
 =:C_{R,t}H_s^2>0.\tag{260}
\]
The corresponding free energy form obeys the previously established total-variation estimate
\[
 \mathcal E_t(J_t h_s)\le \frac{72}{\pi^4t^9}H_s^2.\tag{261}
\]
Consequently
\[
 \frac{\mathcal E_t(J_t h_s)}{\|J_t h_s\|^2}
 \le \frac{72}{\pi^4t^9C_{R,t}}.
 \tag{262}
\]
For fixed $R$, the exact integral in (260) satisfies $C_{R,t}\sim9/(4\pi^4t^8)$ as $t\to\infty$, hence the right side of (262) is asymptotic to $32/t$.  This proves a zero-quotient sequence for the free common-state representation for every nonzero nonnegative compactly supported curvature weight, including the released profile at each preterminal time.  Along the terminal sequence, (254) makes $H_s$ diverge while the same quotient bound remains independent of $H_s$; no norm division was used.

The conclusion still concerns the explicitly identified free comparison Hamiltonian.  Transferring (260)--(262) to the interacting $D_{h_s,g}$ requires the existing positive-coupling norm/energy convergence at each fixed $s,t$ and a jointly chosen $(s,t,g,L)$ trajectory.  The classical source proves neither that joint limit nor a source-free Yang--Mills solution, so this exact free-state zero quotient is a necessary bridge calculation rather than the final mass-gap contradiction.

\section{Diagonal transfer along the proved weak-coupling path}
\label{sec:diagonal-profile-transfer}

The fixed-profile convergence theorem can be combined with the released terminal sequence without normalizing the curvature.  Choose $s_n=1-1/n$ (discarding finitely many indices so that $s_n$ lies in the smooth source interval), set $h_n=\lambda^2|\operatorname{curl}u_\nu(s_n)|^2$, and let $t_n=n$.  Every $h_n$ is nonnegative, smooth and supported in the same compact set.  The support-uniform estimate (260)--(262) gives a positive free norm and a free quotient bounded by $32/n+o(n^{-1})$.

For each $n$, the established positive-coupling theorem supplies a tail of the original sequence
\[
 L_j=j^2,\qquad a_j=(100j)^{-1},\qquad g_j=2^{-m_j}>0,
 \tag{263}
\]
for which the actual centered vectors $y_{h_n,g_j}(t_n)$ have both their raw norm and their energy within any prescribed relative error of the corresponding free quantities.  The profile-dependent error is finite because $h_n$ is smooth and compactly supported at each fixed $n$.  Recursively choose $j_n>j_{n-1}$ in that tail with both relative errors below $1/n$; this is a genuine diagonal subsequence and does not alter $h_n$, $t_n$, $a_j$, or the source units.  Then
\[
 \|y_{h_n,g_{j_n}}(t_n)\|^2
   =\|J_{t_n}h_n\|^2(1+o(1)),\qquad
 \langle y_{h_n,g_{j_n}}(t_n),A_{g_{j_n}}y_{h_n,g_{j_n}}(t_n)\rangle
   =\mathcal E_{t_n}(J_{t_n}h_n)(1+o(1)),
 \tag{264}
\]
and therefore
\[
 \frac{\langle y_{h_n,g_{j_n}}(t_n),A_{g_{j_n}}y_{h_n,g_{j_n}}(t_n)\rangle}
  {\|y_{h_n,g_{j_n}}(t_n)\|^2}
 \le \frac{32}{n}+o(n^{-1})\longrightarrow0.
 \tag{265}
\]
The vectors are vacuum-orthogonal by the exact centering in (258), and no quotient or state norm was used in constructing them.  Equations (263)--(265) are consequently an actual zero-quotient sequence for the changing finite-regulator Hamiltonians on the proved weak-coupling continuum path, with the Navier--Stokes curvature mass retained and diverging along $s_n\uparrow1$.

This does not yet prove failure of the interacting four-dimensional Yang--Mills mass-gap assertion.  The selected couplings satisfy $g_{j_n}\to0$, and the existing convergence theorem controls the free comparison path; it does not prove survival of a fixed renormalized non-Abelian interaction or identify the resulting limit with the target interacting theory.  The diagonal construction isolates the exact remaining requirement: replace (263) by a physical-scale trajectory on which the full interacting Wilson/Haar terms remain present, then repeat the relative-error argument with those terms.  Until that trajectory is proved, (265) is a rigorous weak-coupling counterexample candidate, not a completed Millennium disproof.

\section{Fixed-coupling diagnostic for the competing magnetic channel}
\label{sec:fixed-coupling-diagnostic}

The root Yang--Mills lane supplies an exact fixed-coupling calculation that fixes the next obstruction.  On the unchanged geometric path $L_j=j^2$, $a_j=(100j)^{-1}$, retain $g_j=g_0>0$, so $\xi_0=(4g_0^4)^{-1}$ and $M_{L_j}=12j^4(2j^2+1)$.  For the explicitly defined cusp magnetic translation with $\theta_j=2\pi a_j^2/D_j$, the full interacting vacuum estimate is
\[
 \|\chi_j\|\le \frac{6144\pi^2\xi_0}{10^8}\,j^{-2},
 \qquad
 d_j=\|\chi_j\|^2\le
 \left(\frac{6144\pi^2\xi_0}{10^8}\right)^2j^{-4}.
 \tag{266}
\]
This estimate is obtained from the exact all-angle operator inequality and the full positive vacuum; no constant-vacuum replacement or face deletion is made.  It says that the unnormalized centered measures of this fixed-coupling channel converge to zero in total mass.  Division by $d_j$ leaves a normalized spectral probability, but (266) supplies no low-energy estimate for that probability.

The same dictionary gives $\xi_0M_{L_j}\to\infty$ whenever $g_0$ is fixed and $a_j\downarrow0$.  Therefore the small-interaction regime used by the weak-coupling diagonal is unavailable on this fixed-coupling path.  The exact fixed-coupling result neither proves nor rules out a different non-Abelian state with a vanishing physical quotient; it records why the present NS curvature-state sequence cannot simply be relabelled as such a state.  A fixed-renormalized-coupling counterexample still requires a new state construction whose norm and energy remain controlled when $\xi M_L$ diverges.



